%!TEX root = ../../main.tex

\chapter{Konzeption der experimentellen Evaluierung}
\label{ch:konzeption}

Aufbauend auf den theoretischen Grundlagen aus Kapitel~\ref{ch:grundlagen} und der in Kapitel~\ref{ch:stand-der-forschung} präzisierten Forschungslücke beschreibt dieses Kapitel die Konzeption der experimentellen Evaluierung. Es überführt die in Abschnitt~\ref{subsec:zielsetzung} formulierten Forschungsfragen in ein konkretes Untersuchungsdesign und legt damit die methodische Grundlage für die folgende Umsetzung. Zunächst wird das übergeordnete Untersuchungsdesign und die daraus abgeleiteten Hypothesen dargelegt (Abschnitt~\ref{sec:forschungsdesign}). Anschließend werden die Auswahl der untersuchten \acp{LLM} (Abschnitt~\ref{sec:modellauswahl}), die Definition der betrachteten \ac{ETL}-Aufgaben (Abschnitt~\ref{sec:etl-aufgaben}) sowie die verglichenen Prompting-Strategien (Abschnitt~\ref{sec:prompting-strategien-konzept}) beschrieben. Darauf folgt die Operationalisierung der Evaluationskriterien (Abschnitt~\ref{sec:evaluationskriterien}), bevor abschließend die klassische \ac{ETL}-Pipeline als Vergleichsbasis definiert wird (Abschnitt~\ref{sec:baseline}). Die konkrete technische Umsetzung des hier entworfenen Designs ist Gegenstand der Kapitel~\ref{ch:datensatz} bis~\ref{ch:durchfuehrung}.

\section{Untersuchungsdesign und Hypothesen}
\label{sec:forschungsdesign}

Zur Beantwortung der Forschungsfragen wird ein experimentelles Untersuchungsdesign gewählt, das den Einsatz von \acp{LLM} für ausgewählte \ac{ETL}-Teilaufgaben unter kontrollierten Bedingungen mit einer klassischen, regelbasierten Pipeline vergleicht. Ein experimenteller Zugang ist hier angemessen, da die Forschungslücke gerade im Fehlen belastbarer, vergleichender Evaluierungen unter einheitlichen Bedingungen besteht, während das grundsätzliche Potenzial \ac{KI}-gestützter \ac{ETL}-Prozesse in der Literatur bereits vielfach konzeptionell beschrieben wurde \autocite{walhaDataIntegrationTraditional2024,aAIEnhancedETLProcesses2025}. Den Kern des Designs bildet ein vollständig gekreuztes Versuchsschema über die drei Faktoren Modell, Prompting-Strategie und \ac{ETL}-Aufgabe: Für jede Kombination dieser Faktoren wird die Verarbeitung durchgeführt und anhand einheitlicher, in Abschnitt~\ref{sec:evaluationskriterien} definierter Kriterien bewertet, sodass sich die Einflüsse der einzelnen Faktoren systematisch isolieren lassen.

Eine zentrale Gestaltungsentscheidung betrifft die Einbindung der Modelle in den Verarbeitungsprozess. Als primärer Ausführungsmodus werden sie angewiesen, einen \ac{ETL}-Verarbeitungsschritt als Python-Code zu erzeugen, der anschließend deterministisch auf den Daten ausgeführt und dessen Ergebnis gegen eine bekannte Soll-Lösung (Ground Truth) geprüft wird. Diese Code-Generierung entkoppelt die Verarbeitung von der Datenmenge und damit von der Begrenzung durch das Kontextfenster, erzeugt ein nachvollziehbares und wiederausführbares Artefakt und stellt die unmittelbare Vergleichbarkeit mit der regelbasierten Pipeline her, die ebenfalls programmatisch umgesetzt wird \autocite{khanOverviewETLTechniques2024,nagabhyruBridgingTraditionalETL2022}.

Ergänzend wird ein zweiter Ausführungsmodus untersucht, die direkte In-Context-Verarbeitung, bei der dem Modell die Eingabedaten im Prompt bereitgestellt werden und es die Ergebnistabelle ohne den Umweg über generierten Code zurückgibt. Dieser Modus zielt auf die semantisch geprägten Aufgaben, bei denen sich das erforderliche Wissen nicht vollständig in feste Regeln überführen lässt, etwa bei uneinheitlich geschriebenen Länderangaben oder unscharfen Duplikaten: Das Modell wendet sein Weltwissen unmittelbar auf jeden einzelnen Wert an, statt eine verallgemeinernde Vorschrift zu formulieren \autocite{hassaniRoleChatGPTData2023,khanCognitiveETLUsing2025}. Da die Daten dabei vollständig im Prompt mitgeführt werden, begrenzt das Kontext- und Ausgabefenster diesen Modus auf überschaubare Datenmengen. Er wird deshalb nicht als vollständig gekreuzter Faktor, sondern als ergänzender Vergleich auf denselben Aufgaben und Modellen geführt \autocite{zhaoSurveyLargeLanguage2026,changEfficientPromptingMethods2024}.

Die bis hierhin beschriebene Faktormatrix setzt für jede Aufgabe unmittelbar auf den unveränderten Rohdaten auf. Eine reale \ac{ETL}-Strecke besteht jedoch aus einer Kette, in der jeder Schritt auf dem Ergebnis seines Vorgängers aufsetzt \autocite{khanOverviewETLTechniques2024,walhaDataIntegrationTraditional2024}. Das Design wird daher um eine verkettete Ausführung ergänzt, in der die Ausgabe jedes Schritts die Eingabe des nächsten bildet und ein Abbruch die abhängigen Schritte blockiert, sodass sichtbar wird, wie sich lokal geringe Abweichungen über die Kette hinweg auswirken. Damit ein so gemessener Verlust nicht fälschlich dem betreffenden Schritt zugeschrieben wird, obwohl er aus fehlerhaften Eingabedaten resultiert, tritt ein zweiter, sonst identischer Durchlauf hinzu, in dem jeder Schritt statt der Ausgabe seines Vorgängers die entsprechende Soll-Zwischentabelle erhält. Die Differenz beider Durchläufe beziffert unmittelbar den Anteil des Ergebnisverlusts, der auf vorgelagerte Schritte zurückgeht, eine für FF1 wesentliche Trennung, da eine Empfehlung je Aufgabentyp voraussetzt, dass ein Fehler dem verursachenden Schritt und nicht seinem Nachfolger zugerechnet wird \autocite{nagabhyruBridgingTraditionalETL2022}.

Da \acp{LLM} auch bei identischer Eingabe voneinander abweichende Ausgaben erzeugen können, wird die Reproduzierbarkeit als querschnittliche Dimension in das Design aufgenommen: Jede Konfiguration wird mehrfach ausgeführt, sodass sich die Streuung über die Wiederholungen quantifizieren lässt \autocite{zhaoSurveyLargeLanguage2026,hassaniRoleChatGPTData2023}. Die Wiederholung bezieht sich dabei auf den Generierungsschritt. Das Modell wird jedes Mal erneut zur Code-Erzeugung aufgefordert, während die Ausführung eines einzelnen erzeugten Skripts deterministisch bleibt. Jeder Aufruf erfolgt dabei zustandslos. Aus vorangegangenen Läufen wird nichts übernommen. Steuerbare Einflussgrößen wie die in Abschnitt~\ref{subsec:prompting-einfluss} eingeführte Temperatur bleiben über alle Durchläufe konstant und niedrig, um die nicht modellbedingte Variabilität zu minimieren \autocite{schulhoffPromptReportSystematic2025,zhaoSurveyLargeLanguage2026}.

Aus den Forschungsfragen und dem Forschungsstand lassen sich die folgenden Hypothesen ableiten, die im Rahmen der Auswertung in Kapitel~\ref{ch:ergebnisse} geprüft werden. Zu FF1 werden dabei zwei Hypothesen formuliert, die ihre beiden Teilfragen, den Aufgabentyp und die Ausführungsform, getrennt adressieren, zu FF2 eine. Jede bleibt damit unmittelbar auf eine der beiden Fragen bezogen und bereitet deren Beantwortung vor:

\begin{description}
    \item[H1 (zu FF1):] Die Eignung von \acp{LLM} hängt stark vom Aufgabentyp ab. Bei klar formalisierbaren Aufgaben erreichen sie bestenfalls eine mit der regelbasierten Pipeline vergleichbare Ergebnisqualität, während sie bei Aufgaben mit hoher semantischer Anforderung Vorteile zeigen und mit zunehmender struktureller Komplexität an Zuverlässigkeit verlieren. Ein über alle Aufgaben hinweg überlegener Ansatz existiert daher nicht \autocite{hassaniRoleChatGPTData2023,coteDataCleaningMachine2024}.
    \item[H1a (zu FF1):] Die Eignung wird nicht allein vom Aufgabentyp bestimmt, sondern zusätzlich von der Ausführungsform. Erwartet wird, dass der Ausführungsmodus darüber entscheidet, welche Fähigkeit eines Modells überhaupt zur Wirkung kommt, nämlich Einzelfallverständnis in der Direktverarbeitung und verlässliche Wiederholbarkeit in der Code-Generierung, und dass der Zuschnitt der Aufgabe die Ergebnisqualität stärker beeinflusst als die Wahl des Modells, da mit der Zahl der in einem Arbeitsschritt gleichzeitig zu treffenden Entscheidungen die Fehleranfälligkeit überproportional wächst \autocite{khanCognitiveETLUsing2025,khanOverviewETLTechniques2024}.
    \item[H2 (zu FF2):] Aufwändigere Prompting-Strategien, die zusätzlichen Kontext, Beispiele oder Zwischenschritte bereitstellen, verbessern die Ergebnisqualität gegenüber einem einfachen Zero-Shot-Prompt, wobei der Nutzenzuwachs mit steigendem Aufwand abnimmt und der zusätzliche Kontext bei zu enger Führung auch nachteilig wirken kann \autocite{sahooSystematicSurveyPrompt,schulhoffPromptReportSystematic2025}.
\end{description}

Tabelle~\ref{tab:untersuchungsebenen} fasst die daraus entstehenden Untersuchungsebenen zusammen. Sie unterscheiden sich darin, wie eine Aufgabe gestellt wird, ob unabhängig auf den Rohdaten, als Glied einer Kette oder als Gruppe zusammengefasster Schritte, und welche Eingabe sie dabei erhält. Erst ihr Zusammenspiel erlaubt die Zuordnung eines beobachteten Fehlers zu seiner Ursache: Die Faktormatrix misst die Leistung bei einer isoliert gestellten Aufgabe, die verkettete Ausführung die Leistung im Realbetrieb, und die Differenz zum isolierten Durchlauf beziffert den Anteil, der aus vorgelagerten Schritten stammt.

\begin{table}[H]
    \centering
    \small
    \caption[Untersuchungsebenen des Versuchsdesigns]{Die Ebenen des Untersuchungsdesigns mit ihrer jeweiligen Fragestellung und Eingabeform. Die konkreten Laufzahlen je Ebene nennt Kapitel~\ref{ch:durchfuehrung}.}
    \label{tab:untersuchungsebenen}
    \begin{tabularx}{\textwidth}{@{}lXX@{}}
        \toprule
        \textbf{Ebene} & \textbf{Eingabe je Aufgabe} & \textbf{Beantwortet} \\
        \midrule
        Faktormatrix & Rohdaten, kein Vorgängerergebnis & Einfluss von Modell, Strategie und Aufgabentyp \\
        \addlinespace
        Direkt-Modus & Rohdaten, im Prompt eingebettet & Einfluss des Ausführungsmodus \\
        \addlinespace
        Verkettete Pipeline & Ausgabe des Vorgängerschritts & Leistung im Realbetrieb einer Strecke \\
        \addlinespace
        Isolierter Durchlauf & Soll-Ausgabe des Vorgängers & Trennung eigener Leistung von geerbtem Fehler \\
        \addlinespace
        Gruppierte Strecke & Soll-Ausgabe der Vorgängergruppe & Einfluss des Aufgabenzuschnitts \\
        \addlinespace
        Explizit gegliederte Aufgabe & Rohdaten, benannte Teilschritte & Einfluss der Ausdrücklichkeit der Anforderung \\
        \addlinespace
        Regelbasierte Pipeline & wie die jeweilige Ebene & Vergleichsmaßstab für alle Ebenen \\
        \bottomrule
    \end{tabularx}
\end{table}
\newpage
Die gruppierte Strecke und die explizit gegliederte Aufgabe wurden erst nach einer ersten Sichtung der Ergebnisse ergänzt, um eine im Zuschnitt der Aufgabenstellung erkennbar gewordene Blindstelle zu schließen. Die auf ihnen beruhenden Befunde sind daher als explorativ zu lesen, was bei der Bewertung von Hypothese~H1a zu berücksichtigen ist.

Die Reproduzierbarkeit wird, wie in Abschnitt~\ref{subsec:zielsetzung} festgelegt, nicht als eigenständige Hypothese geführt, sondern als querschnittliche Dimension zu beiden: Für H1 ist die Stabilität Bestandteil der Eignungsbeurteilung, da ein im Mittel gutes, aber stark schwankendes Ergebnis produktiv einen anderen Wert besitzt als ein gleich gutes, verlässlich wiederholbares. Für H2 ist sie zusätzliches Unterscheidungsmerkmal zwischen den Strategien. Gemessen wird, wie stark identische Konfigurationen trotz niedriger Temperatur voneinander abweichen und ob die Stabilität mit zunehmender Aufgabenkomplexität abnimmt \autocite{zhaoSurveyLargeLanguage2026,hassaniRoleChatGPTData2023}.

Der Gegenstand dieser Arbeit ist auf die Transformationsphase des \ac{ETL}-Prozesses beschränkt. Extraktion und Ladevorgang bleiben ebenso außen vor wie Aspekte des produktiven Betriebs, etwa Fehlerbehandlung, Überwachung, inkrementelle Verarbeitung oder Datenschutzanforderungen beim Übertragen von Daten an externe Anbieter. Gerade der letzte Punkt ist für die praktische Anwendbarkeit erheblich und wird hier nur insoweit berührt, als mit dem lokal betriebenen Modell eine Option ohne Datenabfluss mituntersucht wird.

\section{Auswahl der Large Language Models}
\label{sec:modellauswahl}

Die Auswahl der untersuchten Modelle folgt dem Ziel, ein für den praktischen Einsatz repräsentatives und zugleich aussagekräftiges Spektrum aktueller \acp{LLM} abzudecken. Wie in Abschnitt~\ref{subsec:llm-modelle} dargestellt, unterscheiden sich die verfügbaren Modellfamilien nicht nur in Architektur und Parameterzahl, sondern auch in praxisrelevanten Eigenschaften wie Kontextfenster, Verarbeitungsgeschwindigkeit, Nutzungskosten und Verfügbarkeit als offenes oder geschlossenes Modell \autocite{zhaoSurveyLargeLanguage2026,2026AIIndex}. Gerade diese Heterogenität begründet den Vergleich mehrerer Modelle, da sich die Eignung für \ac{ETL}-Aufgaben zwischen ihnen unterscheiden kann und ein Einzelmodell keine verallgemeinerbaren Aussagen zuließe \autocite{ComparisonAIModels,2026AIIndex}.

Als Auswahlkriterien werden die Aktualität und Verbreitung der Modelle, die Abdeckung unterschiedlicher Anbieter sowie die Berücksichtigung sowohl geschlossener als auch offen verfügbarer Modelle herangezogen. Auf dieser Grundlage werden je ein Vertreter der etablierten geschlossenen Modellfamilien von OpenAI, Anthropic und Google sowie ein offen verfügbares, lokal ausführbares Modell aus der Llama-Reihe untersucht. Damit deckt die Auswahl unterschiedliche Kosten-, Leistungs- und Betriebsmodelle ab und erlaubt insbesondere den praktisch bedeutsamen Vergleich zwischen kommerziell betriebenen und lokal betreibbaren Modellen \autocite{zhaoSurveyLargeLanguage2026,2026AIIndex}. Tabelle~\ref{tab:modelle} weist die konkret eingesetzten Modelle mit ihrer Versionskennung und den zugrunde gelegten Preisen aus.

\begin{table}[htbp]
    \centering
    \small
    \caption[Untersuchte Modelle]{Die vier untersuchten Modelle mit Versionskennung und den für die Kostenschätzung zugrunde gelegten Listenpreisen in US-Dollar je einer Million Token (Stand: 27.\,Juli 2026).}
    \label{tab:modelle}
    \begin{tabularx}{\textwidth}{@{}lXlrr@{}}
        \toprule
        \textbf{Anbieter} & \textbf{Versionskennung} & \textbf{Betrieb} & \textbf{Eingabe} & \textbf{Ausgabe} \\
        \midrule
        Anthropic & \texttt{claude-sonnet-5} & Cloud & 2,00 & 10,00 \\
        OpenAI & \texttt{gpt-5.6-terra} & Cloud & 2,50 & 15,00 \\
        Google & \texttt{gemini-3.6-flash} & Cloud & 1,50 & 7,50 \\
        Meta & \texttt{llama3.1:8b} & lokal & \multicolumn{2}{c}{keine nutzungsabhängigen Kosten} \\
        \bottomrule
    \end{tabularx}
\end{table}

\newpage
Bei den drei kommerziellen Modellen handelt es sich jeweils um das mittlere Leistungsniveau der aktuellen Generation des Anbieters, nicht um dessen Spitzenmodell. Alle drei führen intern einen Denkprozess aus, dessen Token zum Ausgabetarif abgerechnet werden und in den ausgewiesenen Kosten enthalten sind. Die angegebenen Preise sind die zum genannten Stichtag herangezogenen Listenpreise.

Um die Modelle technisch austauschbar ansprechen zu können, werden sie hinter einer einheitlichen Schnittstelle gekapselt, sodass der Versuchsablauf unabhängig vom konkreten Anbieter erfolgt. Die Details dieser Abstraktion werden in Kapitel~\ref{ch:implementierung} beschrieben.

\section{Definition der ETL-Aufgaben}
\label{sec:etl-aufgaben}

Die untersuchten \ac{ETL}-Aufgaben orientieren sich an den in Abschnitt~\ref{subsec:dq-probleme} beschriebenen, in der Praxis am häufigsten auftretenden Datenproblemen und decken die drei in dieser Arbeit fokussierten Teilbereiche der Transformationsphase ab: die Datenbereinigung, die Datentransformation und die Duplikaterkennung. Um nicht nur die grundsätzliche Eignung, sondern auch deren Grenzen in Abhängigkeit von der Aufgabenkomplexität zu erfassen, wird jeder Teilbereich in drei aufeinander aufbauenden Schwierigkeitsstufen operationalisiert. Diese reichen von einer einfachen, klar formalisierbaren Aufgabe über eine mittlere Stufe bis hin zu einer anspruchsvollen Variante, die semantisches Verständnis oder mehrstufige Verarbeitung erfordert. Daraus ergibt sich eine Matrix aus neun Aufgaben, die zugleich nach Schwierigkeit gestaffelt ist und so eine differenzierte Beurteilung der Modelle erlaubt \autocite{coteDataCleaningMachine2024,ridzuanReviewDataQuality2024}. Alle Aufgaben beziehen sich auf den synthetischen Datensatz aus Kapitel~\ref{ch:datensatz} und sind so gestaltet, dass für jede von ihnen eine eindeutige Soll-Lösung existiert.

\subsection{Datenbereinigung}
\label{subsec:aufgabe-bereinigung}

Die Aufgaben zur Datenbereinigung adressieren die Beseitigung von Mängeln, die unmittelbar die Genauigkeit, Vollständigkeit und Konsistenz des Datenbestands beeinträchtigen \autocite{coteDataCleaningMachine2024,rangineniReviewEnhancingData2023}. Die einfache Stufe umfasst grundlegende Bereinigungsoperationen wie das Entfernen überflüssiger Leerzeichen in Textspalten und das einheitliche Kennzeichnen fehlender Werte. Diese Aufgabe ist mit wenigen Regeln vollständig beschreibbar und dient als Referenzpunkt für die grundsätzliche Verarbeitungsfähigkeit der Modelle. Die mittlere Stufe verlangt die Vereinheitlichung uneinheitlich formatierter Datumsangaben, die in mehreren unterschiedlichen Darstellungen vorliegen, auf ein einheitliches Zielformat und prüft damit die Fähigkeit zur formatübergreifenden Normalisierung. Die anspruchsvolle Stufe schließlich erfordert die semantische Vereinheitlichung von Länderangaben, die in verschiedenen Sprachen, Abkürzungen und Schreibweisen einschließlich Tippfehlern vorliegen, auf einen standardisierten Code. Diese Aufgabe geht über den rein syntaktischen Abgleich hinaus und setzt ein inhaltliches Verständnis der Werte voraus, das mit starren Regeln nur schwer abzubilden ist \autocite{hassaniRoleChatGPTData2023,khanCognitiveETLUsing2025}.

\subsection{Datentransformation}
\label{subsec:aufgabe-transformation}

Die Aufgaben zur Datentransformation betreffen die Überführung der Daten in die geforderte Struktur und Typisierung sowie die Ableitung neuer Informationen aus vorhandenen Werten. Die einfache Stufe umfasst Typkonvertierungen, etwa die Umwandlung uneinheitlich formatierter Preisangaben in einen numerischen Wert und textueller Wahrheitswerte in einen Boolean-Datentyp, und prüft damit die korrekte Behandlung von Datentypen. Die mittlere Stufe verlangt die Berechnung abgeleiteter Spalten aus bestehenden Feldern, beispielsweise die Ermittlung eines Gesamtbetrags sowie die Extraktion einzelner Datumsbestandteile, und adressiert damit zeilenweise Berechnungslogik. Die anspruchsvolle Stufe bildet die komplexeste Aufgabe der gesamten Untersuchung und verbindet mehrere Verarbeitungsschritte: Sie erfordert die Verknüpfung der drei Tabellen über ihre Schlüsselbeziehungen, eine begleitende Bereinigung der beteiligten Felder sowie eine anschließende Gruppierung und Aggregation der Ergebnisse. Damit prüft sie, ob die Modelle eine mehrstufige, aus Verknüpfung, Bereinigung und Aggregation zusammengesetzte Verarbeitungslogik korrekt umsetzen können, wie sie für reale \ac{ETL}-Strecken charakteristisch ist \autocite{khanOverviewETLTechniques2024,walhaDataIntegrationTraditional2024}.

\subsection{Duplikaterkennung}
\label{subsec:aufgabe-duplikate}

Die Aufgaben zur Duplikaterkennung zielen auf die Eindeutigkeit des Datenbestands und damit auf die Identifikation und Auflösung mehrfach vorhandener Datensätze \autocite{elmagarmidDuplicateRecordDetection2007,coteDataCleaningMachine2024}. Die einfache Stufe beschränkt sich auf exakte Duplikate, bei denen vollständig identische Datensätze zu entfernen sind, und ist mit einem deterministischen Abgleich vollständig lösbar. Die mittlere Stufe verlangt die schlüsselbasierte Entdoppelung anhand eines eindeutigen Merkmals, wobei bei mehreren übereinstimmenden Datensätzen anhand einer definierten Regel der zu behaltende Datensatz ausgewählt werden muss. Die anspruchsvolle Stufe adressiert schließlich unscharfe Duplikate (Fuzzy-Duplikate), bei denen derselbe reale Sachverhalt in abweichenden Schreibweisen vorliegt, etwa durch Tippfehler, unterschiedliche Groß- und Kleinschreibung oder geringfügig abweichende Formatierungen. Da hier ein exakter Abgleich nicht ausreicht und stattdessen eine Ähnlichkeitsbetrachtung im Sinne der von Elmagarmid et al.\ beschriebenen Verfahren erforderlich ist, stellt diese Aufgabe die zweite semantisch geprägte Anforderung dar, an der sich die semantischen Fähigkeiten der Modelle gegenüber regelbasierten Verfahren zeigen \autocite{elmagarmidDuplicateRecordDetection2007,hassaniRoleChatGPTData2023}.

\section{Prompting-Strategien}
\label{sec:prompting-strategien-konzept}

Da \acp{LLM} ihre Aufgaben primär über das In-Context-Lernen lösen, kommt der Gestaltung der Eingabe eine entscheidende Bedeutung für die Ergebnisqualität zu, weshalb die Prompting-Strategie als eigenständiger Untersuchungsfaktor in das Design aufgenommen wird \autocite{sahooSystematicSurveyPrompt,schulhoffPromptReportSystematic2025}. Angesichts der in Abschnitt~\ref{subsec:prompting-strategien} angesprochenen uneinheitlichen Begriffsbildung im Feld werden für diese Arbeit klar abgegrenzte Strategien definiert, die sich in ihrem Aufwand und dem Umfang der bereitgestellten Hilfestellung systematisch unterscheiden und gegeneinander verglichen werden.

Den Ausgangspunkt bildet das Zero-Shot-Prompting, bei dem das Modell die Aufgabe allein anhand einer präzisen Anweisung und ohne weitere Beispiele bearbeitet. Diese Strategie dient als Referenzbasis für die vorliegende Aufgabenbeschreibung \autocite{liuPretrainPromptPredict2023,sahooSystematicSurveyPrompt}. Darauf aufbauend wird das Few-Shot-Prompting untersucht, dessen Wirksamkeit Brown et al.\ an GPT-3 nachgewiesen haben \autocite{brownLanguageModelsFewShot2020}. Es wird hier durch ein einzelnes Beispiel einer anderen Aufgabe operationalisiert und stellt deshalb im engeren Sinne eine One-Shot-Variante dar \autocite{liuPretrainPromptPredict2023}.

Als dritte Strategie wird das Chain-of-Thought-Prompting nach Wei et al.\ einbezogen, bei dem das Modell zur Entwicklung seiner Lösung in nachvollziehbaren Zwischenschritten angeleitet wird, was sich insbesondere bei den mehrstufigen Aufgaben dieser Untersuchung als wirksam erweisen könnte \autocite{weiChainOfThought2022,sahooSystematicSurveyPrompt}. Die Untersuchung beschränkt sich bewusst auf diese drei etablierten Grundformen, um eine reine Aufgabenanweisung, eine beispielgestützte Vorlage und eine explizite Schrittanleitung gegenüberzustellen. Die drei Vorlagen bilden keine vollständige oder streng geordnete Skala des Prompt-Aufwands. Um faire Vergleichsbedingungen zu gewährleisten, werden die Strategien für alle Aufgaben nach einem einheitlichen Aufbau formuliert und versioniert vorgehalten, sodass jede eingesetzte Prompt-Variante eindeutig dokumentiert und wiederverwendbar bleibt \autocite{schulhoffPromptReportSystematic2025,changEfficientPromptingMethods2024}.

\section{Evaluationskriterien}
\label{sec:evaluationskriterien}

Damit die Ergebnisse aller Verfahren über Modelle, Prompting-Strategien und Aufgaben hinweg vergleichbar gemessen werden können, werden die Bewertungskriterien aus den in Abschnitt~\ref{subsec:dq-dimensionen} eingeführten Dimensionen der Datenqualität abgeleitet und operationalisiert \autocite{ridzuanReviewDataQuality2024}. Grundlage jeder Bewertung ist der Abgleich der erzeugten Ausgabedatei mit der für jede Aufgabe bekannten Soll-Lösung. Die folgenden Unterabschnitte beschreiben die vier herangezogenen Kriterien Genauigkeit, Vollständigkeit, Konsistenz und Effizienz. Die Stabilität der Ergebnisse wird, wie in Abschnitt~\ref{sec:forschungsdesign} festgelegt, nicht als separates Einzelkriterium erhoben, sondern daraus abgeleitet, wie stark die genannten Qualitätskennzahlen über die wiederholten Ausführungen einer Konfiguration streuen: Eine geringe Streuung entspricht dabei einer hohen, eine große Streuung einer geringen Stabilität \autocite{zhaoSurveyLargeLanguage2026,hassaniRoleChatGPTData2023}.

\subsection{Genauigkeit}
\label{subsec:genauigkeit}

Die Genauigkeit erfasst, inwieweit die erzeugten Werte mit der Soll-Lösung übereinstimmen, und bildet das zentrale Korrektheitskriterium \autocite{ridzuanReviewDataQuality2024,rangineniReviewEnhancingData2023}. Sie wird als Anteil der korrekt verarbeiteten Werte am Gesamtbestand operationalisiert, sodass auch teilweise korrekte Ergebnisse differenziert bewertet werden. Da programmatisch erzeugte Ergebnistabellen korrekte Werte häufig in abweichender Zeilenreihenfolge liefern, werden beide Tabellen vor dem Abgleich einheitlich sortiert.

Weicht die Datensatzanzahl von der Soll-Lösung ab, lassen sich die Zeilen nicht mehr eindeutig zuordnen. Ein strikter Abgleich würde ein nahezu vollständiges Ergebnis dann genauso bewerten wie ein vollständig falsches und gerade bei der Duplikaterkennung, wo die Datensatzanzahl das eigentliche Ergebnis ist, die gesamte inhaltliche Information verwerfen. Deshalb werden zwei Ausprägungen unterschieden. Die \emph{strenge} Genauigkeit setzt die Übereinstimmung von Spaltenmenge und Datensatzanzahl voraus und beträgt andernfalls null. Die \emph{abgeglichene} Genauigkeit verlangt lediglich die Übereinstimmung des Zielschemas, ordnet die Tabellen über die gemeinsamen Schlüssel beziehungsweise die sortierte Reihenfolge einander zu und verbucht überzählige oder fehlende Datensätze als Abweichung. Numerische Felder werden mit einer Toleranz verglichen.

Als Hauptkennzahl dient die abgeglichene Genauigkeit, da sie den Grad der Aufgabenerfüllung abbildet und nicht nur deren Gelingen oder Scheitern. Die strenge Ausprägung wird herangezogen, wo die exakte Datensatzanzahl selbst Teil der Anforderung ist. Ergänzend weist die \emph{Erfolgsquote} den Anteil der Läufe aus, die eine strukturell verwertbare Ausgabe erzeugt haben, also Zielschema und erwartete Datensatzanzahl einhalten. Ihre Differenz zur Genauigkeit wird nur deskriptiv berichtet, sie lässt keine direkte Bestimmung der Zahl strukturell abweichender Läufe zu. Fehlgeschlagene Ausführungen gehen mit dem Wert null in alle Mittelwerte ein.

Bewertet werden sämtliche Spalten der Soll-Tabelle. Schlüsselspalten werden bei einem schlüsselbasierten Abgleich zur Zuordnung verwendet und anschließend nicht nochmals als Inhaltszellen gezählt, während alle übrigen Soll-Spalten in den Zähler und Nenner eingehen. Bei eindeutigen Schlüsseln werden gleiche Schlüssel über einen Vorkommensindex zugeordnet. Fehlt ein eindeutiger Schlüssel, werden Soll- und Ist-Tabelle nach allen Soll-Spalten als Text sortiert und positionsweise verglichen. Für fehlende oder zusätzliche Zeilen gilt $N=\max(N_{\mathrm{Soll}},N_{\mathrm{Ist}})$, für die bewerteten Inhaltsfelder $C$ damit
\[
    \mathrm{Genauigkeit}=\frac{\text{korrekte Zellen}}{N\cdot C}.
\]
Numerische Werte gelten bei einer relativen Toleranz von $10^{-5}$ und einer absoluten Toleranz von $0{,}01$ als gleich, Textwerte werden nach einer einheitlichen Ersetzung fehlender Werte durch einen Platzhalter exakt verglichen, sodass zwei fehlende Werte als übereinstimmend gelten. Beispiel: Bei einer Soll-Tabelle mit 5 Zeilen, 1 Schlüsselspalte und 4 Inhaltsfeldern sowie einer Ist-Tabelle mit 4 korrekt zugeordneten Zeilen sind $4\cdot4=16$ von $5\cdot4=20$ Zellen korrekt, die abgeglichene Genauigkeit beträgt somit 80~Prozent. Die fehlende Zeile zählt vollständig als Abweichung.

\subsection{Vollständigkeit und Konsistenz}
\label{subsec:vollstaendigkeit}

Die Vollständigkeit gibt an, inwieweit das Ergebnis alle geforderten Datensätze und Felder enthält, ohne dass erwartete Werte fehlen oder unzulässig entfernt werden. Die Konsistenz bezieht sich auf die Widerspruchsfreiheit der Ausgabe, insbesondere auf die Übereinstimmung von Datentypen und Formaten mit dem Zielschema \autocite{ridzuanReviewDataQuality2024,rangineniReviewEnhancingData2023}. Beide werden über den Abgleich mit der Soll-Lösung erfasst, die eine über das Vorhandensein aller erwarteten Spalten und die korrekte Datensatzanzahl, die andere über Schemakonformität und Typisierung. Die Vollständigkeit ist dabei vor allem für die Duplikaterkennung von Bedeutung, bei der sowohl unzureichendes als auch übermäßiges Entfernen sie verletzt, die Konsistenz vor allem für die Transformationsaufgaben, bei denen die korrekte Typisierung ein wesentliches Ergebnismerkmal ist.

Anders als die Genauigkeit werden beide in der Auswertung nicht als eigenständige Kennzahlen berichtet, sondern erfüllen eine vorgelagerte Funktion. Die Vollständigkeit entscheidet darüber, ob ein Ergebnis überhaupt vergleichbar ist, und geht in die abgeglichene Genauigkeit ein, indem fehlende oder überzählige Datensätze dort als Abweichung gewertet werden. Wo der Umfang eigenständige Aussagekraft besitzt, etwa bei der Zahl der entfernten Duplikate, wird er unmittelbar als Datensatzanzahl berichtet. Die Konsistenz wirkt als Ausschlusskriterium: Eine Ausgabe, die das Zielschema verletzt, gilt als strukturell nicht verwertbar, geht mit dem Wert null in die Genauigkeit ein und wird als \emph{Schema-Bruch} gesondert von einem Programmabbruch ausgewiesen. Berichtet werden in Kapitel~\ref{ch:ergebnisse} folglich zwei Qualitätsgrößen, die abgeglichene Genauigkeit und die Erfolgsquote, in denen Vollständigkeit und Konsistenz als strukturelle Voraussetzungen aufgehen, sowie die Effizienzgrößen des folgenden Abschnitts.

\subsection{Effizienz}
\label{subsec:effizienz}

Die Effizienz erfasst den für die Verarbeitung erforderlichen Aufwand und ergänzt die qualitätsbezogenen Kriterien um die für den praktischen Einsatz entscheidende Wirtschaftlichkeitsdimension \autocite{changEfficientPromptingMethods2024,2026AIIndex}. Erhoben werden die Verarbeitungszeit eines Aufrufs, der Verbrauch an Token sowie die daraus resultierenden Nutzungskosten. Diese Kennzahlen sind insbesondere für die vergleichende Bewertung der Prompting-Strategien relevant, da aufwändigere Strategien mit längeren Eingaben einen höheren Token-Verbrauch verursachen und somit zwischen Ergebnisqualität und Effizienz abzuwägen ist \autocite{changEfficientPromptingMethods2024,schulhoffPromptReportSystematic2025}. Zugleich ermöglichen sie den praxisnahen Vergleich zwischen kommerziell betriebenen und lokal ausgeführten Modellen sowie die Gegenüberstellung mit der klassischen Pipeline, deren Ausführung keine vergleichbaren Modellkosten verursacht.

\section{Vergleich gegen eine regelbasierte ETL-Pipeline}
\label{sec:baseline}

Um die Ergebnisse der \ac{LLM}-gestützten Verarbeitung einordnen zu können, wird als Vergleichsmaßstab eine klassische, regelbasiert implementierte \ac{ETL}-Pipeline herangezogen, die dieselben neun Aufgaben auf demselben Datensatz löst. Diese Pipeline bildet eine einfache regelbasierte Referenz, bei dem die Verarbeitungslogik manuell und explizit in Form von Regeln und Skripten vorgegeben wird, und sie wird ebenfalls programmatisch auf Basis der Bibliothek Pandas umgesetzt, um eine technisch vergleichbare Ausgangslage herzustellen \autocite{mishraAutomatingDataIntegration2020}. Sie dient als Referenzpunkt, gegen den sich die Stärken und Schwächen der \ac{KI}-gestützten Ansätze messen lassen, und liefert damit die Grundlage zur Beantwortung von FF1.

Charakteristisch für diese Vergleichsbasis ist, dass sie deterministische und vollständig nachvollziehbare Ergebnisse erzeugt \autocite{walhaDataIntegrationTraditional2024,nagabhyruBridgingTraditionalETL2022}. Damit bildet die regelbasierte Pipeline insbesondere bei den klar formalisierbaren Aufgaben die Referenz, während die semantisch geprägten Aufgaben der Länder-Vereinheitlichung und der Fuzzy-Duplikaterkennung gezielt an die Grenzen rein regelbasierter Verfahren führen. Diese arbeiten zwar auch bei solchen Aufgaben deterministisch, erfassen die zugrunde liegende semantische Ähnlichkeit jedoch nur eingeschränkt. Die Pipeline wird mit den identischen Evaluationskriterien aus Abschnitt~\ref{sec:evaluationskriterien} bewertet, sodass ein unmittelbarer Vergleich zwischen klassischem und \ac{KI}-gestütztem Vorgehen unter gleichen Bedingungen möglich ist und sich auf dieser Grundlage in Kapitel~\ref{ch:handlungsempfehlungen} Handlungsempfehlungen ableiten lassen. Ein hybrider Ansatz, der beide Verfahren kombiniert, ist dabei nicht selbst Gegenstand der Evaluierung, obwohl sich die Nachvollziehbarkeit regelbasierter und die Flexibilität \ac{KI}-gestützter Verfahren in der Praxis häufig ergänzen \autocite{walhaDataIntegrationTraditional2024,kasireddyRoleAIModern2026}.
